\section*{\texorpdfstring{\S\CurrentExerciseGroup}{第{\CurrentExerciseGroup}节}}
\phantomsection
\addcontentsline{toc}{section}{第{\CurrentExerciseGroup}节习题}

\begin{enumerate}
\def\labelenumi{\arabic{enumi})}
\item
  试证：为使数值函数 \(f\)（\(T\) 上的）对 \(T\) 上每个正测度都本质可积，必要且充分的是它有界、有紧支集，且对 \(T\) 上每个正测度都可测。
\item
  设 \(G\) 是桶式空间，\(H\) 是半自反空间，\(F\) 是空间 \(\mathcal{L}_s(G;H)\)。设 \(m\) 是 \(T\) 上的向量测度，取值于 \(F\)。试证：对每个数值函数 \(f\)（关于 \(m\) 本质可积的），有 \(\int fd\mathbf{m} \in F\)。（注意 \(F\) 是半自反的。）\(H = R\) 的情形。
\item
  设 \(\mathbf{m}\) 是 \(\mathrm{T}\) 上的向量测度，取值于 \(\mathrm{F}\)。试证：对每个数值函数 \(f\)（关于 \(\mathbf{m}\) 本质可积的），在下列两种情形的每一种中都有 \(\int f dm \in \mathrm{F}''\)：\(a)\) \(\mathrm{F}\) 是可分辨的（TVS，IV，§2，习题 4）；\(b)\) \(\mathrm{F}\) 是 Fréchet 空间且 \(\mathrm{T}\) 无穷远处可数。（在两种情形中，都把 \(\mathrm{F}\) 嵌入 \(\mathrm{F}''\) 并应用命题 3 的推论 1 的方法；在情形 \(b\) ，用 TVS，IV，§3，第 3 号命题 3 的推论。）
\item
  设 \(\mathbf{m}\) 是 \(\mathrm{T}\) 上的向量测度，取值于 \(\mathrm{F}\)。试证：对每个数值函数 \(f \geqslant 0\)（下半连续且关于 \(\mathbf{m}\) 本质可积的），有 \(\int f d\mathbf{m} \in \mathrm{F}''\)（利用第 IV 章 §1，第 1 号定理 1）。
\item
  设 F 是半自反局部凸空间，m 是 T 上取值于 F 的测度。试证：对每个关于 m 局部可积的数值函数 g，映射 \(f \mapsto \int fg dm\)（\(\mathcal{K}(T)\) 到 F 中的映射）是 T 上的测度，记作 \(g \cdot m\)；为使数值函数 h 关于 \(g \cdot m\) 本质可积，必要且充分的是 gh 关于 m 本质可积，此时有 \(\int h d(g \cdot m) = \int hg dm\)。
\item
  设 F 是半自反局部凸空间，带有一个 R 上的代数结构，使得映射 \((u,v)\mapsto uv\)（\(F\times F\) 到 F 中的映射）对每个变量分别连续。设 m 是 T 上取值于 F 的测度，使得 \(\mathbf{m}(fg)=\mathbf{m}(f)\mathbf{m}(g)\) 对 \(\mathcal{K}(\mathrm{T})\) 中的 f 和 g 成立。试证：若数值函数 f,g 使 f、g 与 fg 都属于 \(\mathcal{L}(\mathbf{m})\)，则 \(\int fg dm=(\int fd\mathbf{m})(\int gd\mathbf{m})\)。（先处理 \(g\in\mathcal{K}(\mathrm{T})\) 的情形，考虑两个测度 \(f\mapsto\mathbf{m}(fg)\) 与 \(f\mapsto\mathbf{m}(f)\mathbf{m}(g)\)；然后类似地处理一般情形并利用习题 5。）\(F=\mathcal{L}(\mathrm{G})\) 的情形，其中 G 自反，F 赋以逐点收敛拓扑或弱逐点收敛拓扑。
\item
  设 \(\mu\) 是 T 上的正测度，m 是 T 上取值于 F 的测度，以 \(\mu\) 为基、以 f 为关于 \(\mu\) 的密度。设 \(q\) 是 F 上的下半连续半范数。试证：若对 T 的每个紧子集 K，\(\int_{K}^{*}(q \circ f) d\mu\) 有限，则对每个 \(\mu\)-可测函数 \(h \geqslant 0\)，不等式 \(\int^{\bullet} h dq(m) \leqslant \int^{\bullet} (q \circ f) h d\mu\) 成立。（利用第 V 章（第 1 版）§2，第 2 号引理 1。）²
\item
  设 \(\mathbf{m}\) 是 \(\mathbf{T}\) 上取值于 \(\mathbf{F}\) 的向量测度，是 \(q\)-可优超的（对一个下半连续半范数 \(q\)，在 \(\mathbf{F}\) 上）。试证：对每个函数 \(f \geqslant 0\)（属于 \(\mathcal{H}(\mathbf{T})\)），
\end{enumerate}

\[
\int f d (q \circ \mathbf {m}) = \sup \sum_ {i} q (\mathbf {m} (f _ {i})),
\]

其中 \((f_i)\) 取遍 \(\mathcal{H}(\mathrm{T})\) 中元素的所有满足 \(\sum_{i}|f_i| \leqslant f\) 的有限序列。（按第 II 章 §2，第 2 号定理 1 的方式论证。）

\begin{enumerate}
\def\labelenumi{\arabic{enumi})}
\setcounter{enumi}{8}
\tightlist
\item
  设 T 是无穷远处可数的局部紧空间，F 是包含可数稠密集的 Hausdorff 桶式空间，m 是 T 上取值于 F 的弱对偶 F' 的向量测度。试证：存在 T 上的正测度 \(\mu\)，使 m 标量意义以 \(\mu\) 为基。（利用第 V 章 §5，第 6 号命题 12，以及 Lebesgue--Nikodym 定理（第 V 章 §5，第 5 号定理 2）推论 5 的条件 5)。）
\end{enumerate}

¶ 10) 设 K 是紧空间。对 K 的每个 Hausdorff 商空间 \(K_{1}\)，把 Banach 空间 \(\mathcal{C}(K_{1})\) 等同于 \(\mathcal{C}(K)\) 的一个闭子空间，借助的是单射 \(f \mapsto f \circ \pi\)，其中 \(\pi\) 是 K 到 \(K_{1}\) 上的典范映射。

\begin{enumerate}
\def\labelenumi{\alph{enumi})}
\item
  设 \(u\) 是 \(\mathcal{C}(\mathrm{K})\) 到拟完备 Hausdorff 局部凸空间 \(\mathrm{F}\) 中的连续线性映射。为使 \(u\) 把 \(\mathcal{C}(\mathrm{K})\) 的单位球变为 \(\mathrm{F}\) 中对 \(\sigma(\mathrm{F},\mathrm{F}')\) 相对紧的子集，只需：对每个可度量化商空间 \(\mathrm{K}_1\)（\(\mathrm{K}\) 的），\(u\) 在 \(\mathcal{C}(\mathrm{K}_1)\) 上的限制具有该性质。（利用 Eberlein 定理（TVS，IV，§5，第 3 号，定理 1），并注意每个序列 \((f_n)\)（在 \(\mathcal{C}(\mathrm{K})\) 中）都含于某个子空间 \(\mathcal{C}(\mathrm{K}_1)\)，其中 \(\mathrm{K}_1\) 是 \(\mathrm{K}\) 的可度量化商空间；为此，考虑连续映射 \(x \mapsto (f_n(x))\)（\(\mathrm{K}\) 到 \(\mathbf{R}^{\mathbf{N}}\) 中的）。）
\item
  对 K 的每个 Hausdorff 商空间 \(K_{1}\)，\(\mathcal{C}(K_{1})\) 到 \(\mathcal{C}(K)\) 中的典范单射的转置是映射 \(\mu\mapsto\pi(\mu)\)（\(\mathcal{M}(K)\) 到 \(\mathcal{M}(K_{1})\) 中的映射）。用 \((\mathcal{M}(\mathrm{K}))'\) 表示 Banach 空间 \(\mathcal{M}(\mathrm{K})\)（即 \(\mathcal{C}(\mathrm{K})\) 的双对偶）的对偶。为使 \(\mathcal{M}(\mathrm{K})\) 的子集 A 对 \(\sigma(\mathcal{M}(\mathrm{K}),(\mathcal{M}(\mathrm{K}))')\) 相对紧，必要且充分的是：对 K 的每个可度量化商空间 \(K_{1}\)，A 在 \(\mathcal{M}(\mathrm{K}_{1})\) 中的典范像对 \(\sigma(\mathcal{M}(\mathrm{K}_{1}),(\mathcal{M}(\mathrm{K}_{1}))')\) 相对紧。（注意 A 是强有界的；于是可设 A 是凸的、均衡的且对淡拓扑 \(\sigma(\mathcal{M}(\mathrm{K}),\mathcal{C}(\mathrm{K}))\) 闭（在 \(\mathcal{M}(\mathrm{K}_{1})\) 中应用第 IV 章 §7 习题 10）；因此 A 是对偶 \(F'\)（某个 Banach 空间 F 的对偶）的单位球在一个连续线性映射 u（\(\mathcal{C}(\mathrm{K})\) 到 F 中的）的转置下的像（考虑 A 在 \(\mathcal{C}(\mathrm{K})\) 中的极）；然后应用 a）。）
\end{enumerate}

¶ 11) 设 F、G 是两个 Hausdorff 局部凸空间，其中 G 假设为拟完备的；设 u 是 F 到 G 中的连续线性映射；其转置 \(^{t}u\) 是强连续的，故转置 \(^{t}(^{t}u)\) 是 F'\,' 到 G'\,' 中的强连续且弱连续的线性映射。试证下列条件等价：

\(\alpha)\) \(u\) 把 \(\mathbf{F}\) 的每个有界子集变为 \(\mathbf{G}\) 中对 \(\sigma(\mathbf{G},\mathbf{G}')\) 相对紧的子集。

\(\beta)^{t}(^{t}u)\) 把 \(\mathbf{F}^{\prime \prime}\) 映到 G 中。

\(\gamma)\ ^{t}u\) 对拓扑 \(\sigma(\mathrm{G}^{\prime},\mathrm{G})\) 与 \(\sigma(\mathrm{F}^{\prime},\mathrm{F}^{\prime\prime})\) 连续。

\(\delta)^{t}u\) 把 \(\mathrm{G}'\) 的每个等度连续子集变为对 \(\sigma (\mathbf{F}',\mathbf{F}'')\) 相对紧的集。

（利用下述事实：F 的每个有界子集在 \(F''\) 中对 \(\sigma(F'', F')\) 相对紧。为看出 \(\delta\) 蕴含 \(\gamma\)，先考虑 G 完备的情形，并用 TVS，III，§3，第 6 号定理 2 的推论 1。为过渡到一般情形，把 G 嵌入其完备化空间，并注意 \(F''\) 的每个点都在 F 的某个有界子集的闭包中（对拓扑 \(\sigma(F'', F')\)）。）

¶ 12) a) 设 K 是紧空间，E = E(K) 是 R \(^{K}\) 的由开集特征函数的线性组合所构成的子空间；把 E 等同于 C(K) 的双对偶的一个子空间。试证：对 σ(M(K),E) 的 Cauchy 序列对 σ(M(K),(M(K))') 收敛（利用第 V 章 §5 的命题 12 与习题 15）。由此推出：M(K) 的每个对拓扑 σ(M(K),E) 相对紧的子集 A，都对拓扑 σ(M(K),(M(K))') 相对紧。（利用习题 10 b)，化归到 K 可度量化的情形。应用第 V 章 §5 习题 13 证明 A 有界。注意在 A 上，淡拓扑 σ(M(K),C(K)) 与拓扑 σ(M(K), \(\overline{E}\)) 相同，其中 \(\overline{E} \supset C(K)\) 是 E 在 M(K) 的强对偶中的闭包；由此并利用 C(K) 可分这一事实，推出 A 的点的每个序列都有对 σ(M(K),E) 的 Cauchy 子序列；最后援引 Eberlein 定理完成证明。）

\begin{enumerate}
\def\labelenumi{\alph{enumi})}
\setcounter{enumi}{1}
\item
  设 T 是局部紧空间，m 是 T 上取值于拟完备 Hausdorff 局部凸空间 F 的向量测度。假设对 T 的每个紧子集 K，\(\int_{K} dm \in F\)；试证：对 T 上每个有紧支集的有界 Borel 函数 f，有 \(\int f dm \in F\)，且在映射 \(f \mapsto \int f dm\) 之下，支集含于 T 的紧子集 K 且范数 \(\|f\| \leqslant 1\) 的 Borel 函数所成之集的像对 \(\sigma(F, F')\) 相对紧（利用 a) 与习题 11，把它应用于 \(u : f \mapsto \int f dm\)，其中 f 取遍 T 的紧子集的特征函数的线性组合所成之集）。
\item
  设 F 是拟完备的 Hausdorff 局部凸空间，m 是 T 上取值于 F 的向量测度，标量意义以 \(\mu\) 为基；对每个 \(z^{\prime} \in F^{\prime}\)，记 \(z^{\prime} \circ m = g_{z^{\prime}} \cdot \mu\)。为使 b) 的条件成立，必要且充分的是：对 T 的每个紧子集 K、每个等度连续子集 \(H^{\prime}\)（\(F^{\prime}\) 的）以及每个 \(\varepsilon > 0\)，存在 \(\delta > 0\)，使得关系式 \(A \subset K\) 与 \(\mu^{*}(A) \leqslant \delta\) 蕴含 \(\int_{A}^{*}|g_{z^{\prime}}| d\mu \leqslant \varepsilon\)，对每个 \(z^{\prime} \in H^{\prime}\) 成立（用上面的习题 11 与第 V 章 §5 习题 15）；这时称 m 关于 \(\mu\) 绝对连续（对 F 的原拓扑）。每个取值于半自反空间且标量意义以 \(\mu\) 为基的向量测度都关于 \(\mu\) 绝对连续（命题 3 的推论 2）。每个取值于 Banach 空间 \(\mathbf{F}\) 的可优超向量测度 m 都关于 \(|\mathbf{m}|\) 绝对连续（对 \(\mathbf{F}\) 的原拓扑）。
\end{enumerate}

¶ 13) 设 G 是 Hausdorff 桶式空间，F = G' 是其对偶；若当 F 赋以一个与 F 和 G 之间的对偶相容的拓扑时，m 是 T 上取值于 F 的向量测度，则当 F 赋以 G 的对偶的强拓扑时，m 仍是向量测度。

假设 m 标量意义以 \(\mu\) 为基；则 m 关于 \(\mu\) 绝对连续，当 F 赋以拓扑 \(\tau(\mathrm{F},\mathrm{G})\) 时（习题 12 c)）。为使当 F 赋以强拓扑时 m 关于 \(\mu\) 绝对连续，必要且充分的是：对 T 的每个紧子集 K 以及每个序列 \((\mathrm{A}_{n})\)（由 K 的 \(\mu\)-可测子集组成），在 \(\mathcal{L}^{1}(\mu)\) 中，由绝对值 \(\leqslant1\)、且在由序列 \(A_{n}\) 生成的集环（第 IV 章 §4，第 9 号）上的阶梯函数所成之集的闭包在 m 下的像 B，都包含一个对 \(G'\) 的强拓扑稠密的可数子集。（为看出条件是充分的，化归到 T = K 是 Stone 空间的情形（第 IV 章 §4，习题 10）。然后用反证法论证：若 \((\mathrm{A}_{n})\) 是由 K 的既开又闭的子集组成的序列，则在过渡到 K 的一个可度量化商空间 \(K_{1}\)（习题 10 a)）之后，可以假设 \(K_{1}\) 上与 \(\varphi_{A_{n}}\) 相对应的连续函数在 \(\mathcal{C}(\mathrm{K}_{1})\) 中组成一个完全集。另一方面，考虑 G 对与 B 正交的子空间的商，可以假设 B 在 \(F = G'\) 中对拓扑 \(\sigma(\mathrm{G}',\mathrm{G})\) 是完全的；若 H 是由 B 生成的 \(F = G'\) 的子空间，则假设蕴含 G 的每个有界子集 C 都对 \(\sigma(\mathrm{G},\mathrm{H})\) 预紧且可度量化。因此 G 的点的每个有界序列 \((\mathbf{z}_{n})\) 都有一个对 \(\sigma(\mathrm{G},\mathrm{H})\) 的 Cauchy 子序列；试证：若 \((\mathbf{z}_{n})\) 是对 \(\sigma(\mathrm{G},\mathrm{H})\) 的 Cauchy 序列，且记 \(z_{n} \circ m = g_{z_{n}} \cdot \mu\)，则 \(g_{z_{n}}\) 的类所成的序列在 \(L^{1}(\mu)\) 中对拓扑 \(\sigma(\mathrm{L}^{1},\mathrm{L}^{\infty})\) 是 Cauchy 序列；最后，利用第 V 章 §5 的习题 15 与 16 由此得出矛盾。）

\begin{enumerate}
\def\labelenumi{\arabic{enumi})}
\setcounter{enumi}{13}
\tightlist
\item
  \begin{enumerate}
  \def\labelenumii{\alph{enumii})}
  \tightlist
  \item
    设 G 是 Banach 空间，\(F = G'\) 是其强对偶，g 是 T 到 F 中的映射，属于空间 \(\Lambda_{G'}^{p}\)（§1，习题 16）。试证：若 p \textgreater{} 1，则测度 \(g \cdot \mu\) 关于 \(\mu\) 绝对连续（对 F 的强拓扑）。
  \end{enumerate}
\end{enumerate}

\begin{enumerate}
\def\labelenumi{\alph{enumi})}
\setcounter{enumi}{1}
\tightlist
\item
  取 G 为 Banach 空间 \(\mathrm{L}^{1}(\mathbf{N})\)，于是 \(\mathrm{G}^{\prime}=\mathrm{L}^{\infty}(\mathbf{N})\)；若 f 是 §1 习题 7a) 中定义的函数（视为取值于 \(G^{\prime}\) 中），则 f 是标量良可积的，但测度 \(f\cdot\mu\) 不关于 \(\mu\) 绝对连续（对 \(G^{\prime}\) 的强拓扑）。
\end{enumerate}

¶ 15) a) 设 f 是 T 到拟完备 Hausdorff 局部凸空间 F 中的标量局部可积映射。则 \(m = f \cdot \mu\) 是取值于 \(F'^{*}\) 的向量测度（对拓扑 \(\sigma(F'^{*}, F')\)）。试证：若这一测度是绝对连续的（对在 \(F'\) 的等度连续子集上一致收敛的拓扑），且此外 f 是 \(\mu\)-可测的（对 F 的原拓扑），则 m 是取值于 F 的向量测度。（利用 §1 第 2 号命题 8。）若 F 是 Banach 空间，且对某个 p \textgreater{} 1，\(\langle z', f \rangle\) 属于 \(\overline{\mathcal{L}}^{p}(\mu)\)（对一切 \(z' \in F'\)），则测度 m 是绝对连续的（参见习题 14 a））。

\begin{enumerate}
\def\labelenumi{\alph{enumi})}
\setcounter{enumi}{1}
\tightlist
\item
  设 \(\mathrm{I} = [0,1]\)，\(f\) 是定义在 \(\mathrm{I} \times \mathrm{I}\) 上（利用连续统假设）的函数，使得对每个 \(t \in \mathrm{I}\)，\(s \mapsto f(s,t)\) 是一个可数集的特征函数，而对每个 \(s \in \mathrm{I}\)，\(t \mapsto f(s,t)\) 是一个可数集的余集的特征函数（第 V 章 §8，习题 7 c)）。用 \(\mathrm{I}_0\) 表示赋有离散拓扑的区间 I，用 F 表示 I 上有界函数的 Banach 空间 \(\mathcal{B}(\mathrm{I}) = \mathrm{L}^\infty(\mathrm{I})\)；可把 F 与空间 \(\mathcal{C}(\widetilde{\mathrm{I}}_0)\)——即 \(\mathrm{I}_0\) 的 Stone-Čech 紧化上的连续函数所成的空间——等同起来。对每个 \(s \in \mathrm{I}\)，用 \(\mathbf{g}(s)\) 表示 F 的元素 \(t \mapsto f(s,t)\)；试证 \(\mathbf{g}\) 对 I 上的 Lebesgue 测度 \(\mu\) 是标量可积的；但 \(\int \mathbf{g} d\mu\) 不属于 F。更确切地说，\(\mathbf{g} \cdot \mu\) 是取值于 \(\mathrm{F}''\) 的向量测度，等于 \(\mathbf{c}\mu\)，其中 \(\mathbf{c}\) 是与 \(\widetilde{\mathrm{I}}_0 - \mathrm{I}_0 = \mathrm{E}\) 的特征函数等同的常向量（参见习题 12 a））。（把 \(\widetilde{\mathrm{I}}_0\) 上的每个正测度分解为一个原子测度与一个弥散测度之和（第 V 章 §5，第 10 号，命题 15），并注意弥散测度由 E 承载。）由此推出：不存在任何 \(\mu\)-可测函数 \(\mathbf{h}\)（取值于 \(\mathbf{F}\) 的），使 \(\mathbf{h} - \mathbf{g}\) 是标量 \(\mu\)-可忽略的（利用 \(a\)））。
\end{enumerate}

\begin{enumerate}
\def\labelenumi{\arabic{enumi})}
\setcounter{enumi}{15}
\item
  考察映射 \(u_{m}\)——T = {[}0,1{]} 到一个可分 Hilbert 空间中的、§1 习题 1 中定义的映射——并令 \(f_{k} = u_{1} + \cdots + u_{k}\)；试证向量测度 \(f_{k} \cdot \mu\) 在 \(\mathcal{C}(T)\) 的每个有界子集上一致收敛于一个取值于 F 的向量测度 m。此外，m 可连续延拓到空间 \(\mathcal{L}^{2}(\mu)\)，且对每个函数 \(f \in L^{2}\)，有 \(\int f dm \in F\) 以及在 F 中 \(|\int f dm| \leqslant N_{2}(f)\)；特别地，m 标量意义以 \(\mu\) 为基，且对强拓扑关于 \(\mu\) 绝对连续（习题 12 c)）。然而，对 T 上任何正测度 \(\nu\)，m 都不以 \(\nu\) 为基，从而（第 5 号，定理 1 的推论 4）更不是可优超的。（利用第 V 章 §5，第 7 号定理 3，化归到 \(\nu\) 以 \(\mu\) 为基的情形。）
\item
    设 \(\mu\) 是 \(T = [0,1]\) 上的 Lebesgue 测度；映射 \(f \mapsto f \cdot \mu\)（\(\mathcal{C}(T)\) 到 \(\mathcal{M}(T)\) 中的）对 \(\mathcal{M}(T)\) 上的强拓扑连续，因而是取值于该 Banach 空间的向量测度 \(\mathbf{m}\)。试证 \(\mathbf{m}\) 标量意义以 \(\mu\) 为基，且对强拓扑关于 \(\mu\) 绝对连续，但 \(\mathbf{m}\) 不以 \(\mu\) 为基（对 \(\mathcal{M}(T)\) 上的强拓扑）。（注意 \(\mathbf{m}\) 以 \(\mu\) 为基（对淡拓扑 \(\sigma(\mathcal{M}(T), \mathcal{C}(T))\) 而言），并由此推出：若有 \(\mathbf{m} = \mathbf{g} \cdot \mu\)，其中 \(\mathbf{g}\) 标量 \(\mu\)-可积（对 \(\mathcal{M}(T)\) 上的强拓扑），则几乎处处必有 \(\mathbf{g}(t) = \varepsilon_t\)。试证这导致矛盾，办法是注意：对每个有界数值函数 \(\theta\)（\(T\) 上的），线性形式 \(\mathbf{z}' : \lambda \mapsto \sum_{t \in T} \theta(t) \lambda(\{t\})\)
\end{enumerate}

在 \(\mathcal{M}(\mathrm{T})\) 上连续，但 \(\langle \mathbf{g},\mathbf{z}'\rangle\) 不一定 \(\mu\)-可测。）

¶ 18) a) 设 T 是局部紧空间，\((K_{\alpha})\) 是 T 的由紧集组成的局部有限覆盖。设 \(\mu\) 是 T 上的正测度，\(\mu_{\alpha}\) 是 \(\mu\) 在 \(K_{\alpha}\) 上诱导的测度。假设每个空间 \(L^{\infty}(K_{\alpha}, \mu_{\alpha})\) 都具有提升性质。试证 \(L^{\infty}(T, \mu)\) 具有提升性质。

\begin{enumerate}
\def\labelenumi{\alph{enumi})}
\setcounter{enumi}{1}
\tightlist
\item
  设 K 是可度量化紧空间，\(\mu\) 是 K 上的正测度，\((\varpi_n)\) 是由 K 的分成可积集的有限分割组成的基本序列（第 IV 章 §5，习题 17）。对 K 的每个可积子集 A 以及每个元素 \(\widetilde{f} \in \mathrm{L}^1(\mu)\)，令 \(\lambda_A(\widetilde{f}) = 0\)（若 \(\mu(A) = 0\)），否则令 \(\lambda_A(\widetilde{f}) = (\mu(A))^{-1} \int_{A} f d\mu\)；对每个 \(n\)，令 \(\rho_n(\widetilde{f}) = \sum_k \lambda_{A_k}(\widetilde{f}) \varphi_{A_k}\)，如果
\end{enumerate}

\(\varpi_{n} = (\mathrm{A}_{k})\)。试证序列 \((\rho_{n}(\tilde{f}))\)（\(\mu\)-可积函数的序列）几乎处处收敛于 \(f\)。（先考虑 \(f\) 是连续函数的情形。给定任意数 \(a > 0\)，再证明：所有属于至少一个分割 \(\varpi_{n}\) 且满足 \(a \cdot \mu(A) \leqslant \int_{A} f d\mu\) 的集 A 的并 B 是可测的，且 \(\mu(B) \leqslant a^{-1} \int f d\mu\)。最后对 \(f\)，用一列连续函数在 \(\mathcal{L}^{1}\) 中逼近它。）

\begin{enumerate}
\def\labelenumi{\alph{enumi})}
\setcounter{enumi}{2}
\item
  试证每个可度量化紧空间都具有提升性质。（设 \(\mathfrak{U}\) 是 \(\mathbf{N}\) 上比 Fréchet 滤子更细的超滤子；试证 \(\rho(\widetilde{f}) = \lim_{\mathfrak{U}} \rho_n(\widetilde{f})\) 是 \(\mathrm{L}^{\infty}\) 的一个提升。）
\item
  由 \(a\) 与 \(c\) 推出：每个可度量化局部紧空间都具有提升性质（对任何正测度均是）。
\end{enumerate}

\begin{enumerate}
\def\labelenumi{\arabic{enumi})}
\setcounter{enumi}{18}
\tightlist
\item
  设 \(\mu\) 是 T 上的测度，使得 Banach 空间 \(L^1 (\mu)\) 可分。试证：\(L^1 (\mu)\) 到强对偶 \(F'\)（任意 Banach 空间 \(F\) 的强对偶）中的每个连续线性映射，都可以由形如 \(g\mapsto \int g\mathbf{f}d\mu\) 的映射通过过渡到商而得到，其中 \(\mathbf{f}\in \mathcal{L}_{\mathrm{F}_s'}^\infty\)。（利用 TVS，IV，§2 习题 19，化归到 Dunford-Pettis 定理。）
\end{enumerate}

¶ 20) 设 F 是可分 Fréchet 空间，F' 是其对偶，μ 是 T 上的正测度，m 是 T 上取值于弱对偶 \(F'_s\) 的测度，标量意义以 μ 为基。为使 m 以 \(\mu\) 为基，必要且充分的是下列条件成立：对每个 \(\varepsilon > 0\) 以及每个紧子集 \(K\)（\(T\) 的），存在紧子集 \(K_1 \subset K\)，使 \(\mu(K - K_1) \leqslant \varepsilon\)，且使得在 \(m\) 之下，由所有 \(\mu\)-可测函数 \(g\)（有界、支集含于 \(K_1\) 且满足 \(\int |g| d\mu \leqslant 1\) 的）组成的集的像，是 \(F'\) 的等度连续子集。（回忆：\(\int g dm \in F'\) 对每个有界且具有紧支集的 \(\mu\)-可测函数 \(g\) 都成立；参见命题 3 的推论 2。为证条件必要，用 §1 第 5 号命题 13 与 §1 第 2 号命题 5。为证条件充分，先考虑 \(T\) 紧的情形；应用第 5 号定理 1 的推论 3，先证明存在 \(T\) 分解为一个可忽略集 \(N\) 与一列紧集（\(K_n\)）的分割，以及一个可测映射 \(g\)（\(T\) 到 \(F'_s\) 中的），使得 \(\int f dm = \int g f d\mu\) 对每个函数 \(f \in \mathcal{L}^\infty(\mu)\) 都成立，只要该函数在有限个集 \(K_n\) 的并的余集上为零；为证明 \(m = g \cdot \mu\)，用第 V 章 §5 的习题 16 b) 与 20 a)。最后，为过渡到 \(T\) 是任意局部紧空间的情形，用第 IV 章 §5，第 9 号命题 14。）

¶ 21) 设 F 是 Banach 空间，u 是 Banach 空间 \(\mathrm{L}_{\mathrm{F}}^{p}(\mathrm{T},\mu)\) 上的连续线性形式；用 q 表示 p 的共轭指数，用 E 表示 \(\mu\)-可积集上的数值阶梯函数所成的空间。对 \(f \in L^{p}\) 与 \(z \in F\)，可写 \(u(fz) = \langle z, m(f) \rangle\)，其中 m 是 \(L^{p}\) 到 \(F'\) 中的连续线性映射，满足 \(|m(f)| \leqslant \|u\| \cdot N_{p}(f)\)。

\begin{enumerate}
\def\labelenumi{\alph{enumi})}
\tightlist
\item
  设 \((\mathrm{A}_i)_{1\leqslant i\leqslant n}\) 是由两两不交且非可忽略的 \(\mu\)-可积集组成的有限序列。试证
\end{enumerate}

\[
\sum_ {i} \left(| \mathbf {m} (\varphi_ {\mathrm{A} _ {i}}) | ^ {q} / (\mu (\mathrm{A} _ {i})) ^ {q - 1}\right) \leqslant \| u \| ^ {q}.
\]

（对每个向量系统 \((\mathbf{a}_i)_{1\leqslant i\leqslant n}\)（F 中的向量，满足 \(|\mathbf{a}_i| = 1\)），考虑线性形式 \((\xi_i)\mapsto u\big(\sum_i\xi_i\mathbf{a}_i\varphi_{\mathrm{A}_i}\big)\)（\(\mathbf{R}^n\) 上的），并对一个有限支集的测度应用第 V 章 §5，第 8 号定理 4。）由此推出：若 \(\mathrm{A} = \bigcup_{i}\mathrm{A}_{i}\)，则 \(\sum_{i}|\mathbf{m}(\varphi_{\mathrm{A}_i})|\leqslant \| u\| (\mu (\mathrm{A}))^{1 / p}\)（应用 Minkowski 不等式）。

\begin{enumerate}
\def\labelenumi{\alph{enumi})}
\setcounter{enumi}{1}
\tightlist
\item
  对每个正函数 \(f \in \mathcal{E}\)，令 \(\nu(f) = \sup\left(\sum_{i} |\mathbf{m}(f\varphi_{\mathrm{A}_i})|\right)\)，其中 \((\mathrm{A}_i)\)
\end{enumerate}

取遍由两两不交的 \(\mu\)-可积集组成的有限序列所成之集。试证 \(\nu\) 是某个正测度在 \(\mathcal{E}\) 上的限制——该正测度以 \(\mu\) 为基（仍记作 \(\nu\)），定义在 T 上——使得 \(|\mathbf{m}(f)| \leqslant \nu(f)\) 对每个正函数 \(f \in \mathcal{E}\) 都成立（利用 \(a\)））。

\begin{enumerate}
\def\labelenumi{\alph{enumi})}
\setcounter{enumi}{2}
\tightlist
\item
  试证：若 F 可分，则存在 T 到 F' 中的映射 g，标量 μ-可测，使得 \(u(\widetilde{\mathbf{f}})=\int\langle\mathbf{f},\mathbf{g}\rangle d\mu\) 对每个函数 \(f\in L_{F}^{p}\) 都成立，且 \(\|u\|=N_{q}(\mathbf{g})\)。（为证最后的等式，用 §1 第 5 号命题 13 与 §1 习题 13。）d) 由 c) 推出：若 F 是自反可分 Banach 空间，则空间 \(L_{F}^{p}\) 对 \(1<p<+\infty\) 自反（应用 §1 第 5 号命题 12）。
\end{enumerate}

¶ 22) 设 F 是 Hausdorff 局部凸空间，S 是 F 的子集所成之集，这些子集是凸的、均衡的且对 \(\sigma(F, F')\) 紧，而 \(S'\) 是 \(F'\) 的子集所成之集，这些子集是凸的、等度连续的、均衡的且对 \(\sigma(F', F'')\) 紧。

\begin{enumerate}
\def\labelenumi{\alph{enumi})}
\item
  假设 \(\mathfrak{S}\) 中每个集都对 \(\mathfrak{S}'\)-拓扑预紧。试证：每个连续线性映射 \(u\)——\(F\) 到拟完备 Hausdorff 局部凸空间 \(G\) 中的连续线性映射——只要把 \(F\) 的有界子集变为对 \(\sigma(G, G')\) 相对紧的子集，就把每个属于 \(\mathfrak{S}\) 的集变为 \(G\) 中对原拓扑相对紧的子集。（用 TVS，IV，§1 习题 4 与上面的习题 11。）逆命题（考虑 \(F\) 到其对 \(\mathfrak{S}'\)-拓扑的完备化空间中的典范映射）。
\item
  假设 F 是可度量化空间，或是可度量化空间的严格直接极限。试证：若每个对 \(\sigma(\mathrm{F},\mathrm{F}^{\prime})\) 的 Cauchy 序列都是对 \(S^{\prime}\)-拓扑的 Cauchy 序列，则 a) 的条件满足（利用 Šmulian 定理（TVS，IV，§5，习题 2 c）））。
\item
  假设 F 是拟桶式的（TVS，III，§4，习题 7），并用 \(\mathfrak{S}''\) 表示 \(\mathrm{F}''\) 的子集所成之集，这些子集是凸的、等度连续的、均衡的且对 \(\sigma(\mathrm{F}''\) , \(\mathrm{F}'''\) ) 紧。试证：若 \(\mathfrak{S}'\) 中每个集都对 \(\mathfrak{S}''\)-拓扑预紧，则 \(\mathfrak{S}\) 中每个集都对 \(\mathfrak{S}'\)-拓扑预紧。（用 TVS，IV，§1 习题 4，并注意 F 的原拓扑由 \(\mathrm{F}''\) 上的强拓扑诱导。）
\end{enumerate}

¶ 23) 设 T 是局部紧空间，\(\mu\) 是 T 上的正测度，F 是三个 Banach 空间 \(\overline{\mathcal{K}(T)}\)、\(L^1(\mu)\)、\(L^\infty(\mu)\) 之一。设 u 是 F 到拟完备空间 G 中的连续线性映射，它把 F 的单位球变为 G 中对 \(\sigma(G, G')\) 相对紧的子集。试证 u 把 F 的每个对 \(\sigma(F, F')\) 相对紧的子集变为 G 中对原拓扑相对紧的子集。（应用习题 22 c)，把情形 \(F = L^1(\mu)\) 化归到情形 \(F = L^\infty(\mu)\)；用第 II 章 §1 习题 13 把情形 \(L^\infty(\mu)\) 化归到情形 \(F = \mathcal{C}(S)\)，其中 S 是一个适当的紧空间。若 \(F = \overline{\mathcal{K}(T)}\)，则应用习题 22 b)，然后用第 V 章 §5 的习题 22 b) 与 15；注意对 \(\sigma(F, F')\) 的 Cauchy 序列在 T 上逐点收敛，从而对 T 上每个有界测度依测度收敛。）

¶ 24) 设 \(\mu\) 是 \(T\) 上的正测度，\(u\) 是 \(L^1(\mu)\) 到 Banach 空间 \(F\) 中的线性映射，把 \(L^1(\mu)\) 的单位球变为 \(F\) 中对 \(\sigma(F, F')\) 相对紧的子集。试证：存在 \(\mu\)-可测映射 \(f\)（\(T\) 到 \(F\) 中的），使 \(|\mathbf{f}(t)| \leqslant \|u\|\) 对一切 \(t \in T\) 成立，且使得

\[
u (\widetilde {g}) = \int \mathbf {f} g d \mu
\]

对每个函数 \(g \in \overline{\mathcal{L}}^1(\mu)\)（Dunford-Pettis-Phillips 定理）。（先借助第 IV 章 §5，第 9 号命题 14 化归到 T 紧的情形。此时 \(L^\infty(\mu) \subset L^1(\mu)\) 的单位球 B 对 \(\sigma(L^1, L^\infty)\) 相对紧（第 V 章 §5，习题 15）；利用习题 23，试证由 \(u(L^1)\) 生成的 F 的闭线性子空间是可分的。于是可化归到 F 可分且 \(\overline{u(L^1)} = F\) 的情形。此时在 \(u\) 之下，\(L^1\) 的单位球的像在 F 中的闭包 A 对 \(\sigma(F, F')\) 紧（第 IV 章 §7，习题 10）；试证它对该拓扑可度量化（TVS，III，§3，第 4 号，命题 6 的推论 2）。注意此时 A 可等同于一个赋范空间 G 的对偶的单位球，其中 G 是把 \(F'\) 赋以等于 \(A^\circ\) 的规范函数的范数而得到的。试证 G 可分，从而化归到应用 Dunford-Pettis 定理（定理 1 的推论 2）。

¶ 25) 设 \(\mu\) 是 T 上的正测度。

\begin{enumerate}
\def\labelenumi{\alph{enumi})}
\item
  设 \(\mathbf{f}\) 是标量 \(\mu\)-可测映射（\(\mathrm{T}\) 到 Banach 空间 \(\mathrm{F}\) 中的），使得 \(\mathbf{f}(\mathrm{T})\) 对 \(\sigma(\mathrm{F},\mathrm{F}')\) 相对紧。试证：存在 \(\mu\)-可测映射 \(\mathbf{g}\)（\(\mathrm{T}\) 到 \(\mathrm{F}\) 中的），使 \(\mathbf{f} - \mathbf{g}\) 是标量局部可忽略的。（考虑映射 \(\widetilde{h} \mapsto \int \mathbf{f}h d\mu\)（\(\mathrm{L}^1(\mu)\) 到 \(\mathrm{F}\) 中的），对它应用习题 24，并用第 IV 章 §7 习题 10。）
\item
  设 \(\mathbf{f}\) 是 T 到自反 Fréchet 空间 F 中的标量 \(\mu\)-可测映射。试证：存在 T 到 F 中的 \(\mu\)-可测映射 \(\mathbf{g}\)，使 \(\mathbf{f} - \mathbf{g}\) 是标量局部可忽略的（参见习题 15 b））。（用第 IV 章 §5，第 9 号命题 14，化归到 T 紧的情形。然后应用 §1 习题 23 c)，化归到 \(\mathbf{f}(T)\) 在 F 中有界的情形。最后把 F 嵌入可数个 Banach 空间的乘积，并用 a)。）
\item
  设 \(\mathbf{f}\) 是 \(\mathrm{T}\) 到 Fréchet 空间 \(\mathrm{F}\) 中的映射，它是 \(\mu\)-可测的（对拓扑 \(\sigma(\mathrm{F}, \mathrm{F}')\)）。试证 \(\mathbf{f}\) 是 \(\mu\)-可测的（对 \(\mathrm{F}\) 的原拓扑）。（化归到 \(\mathrm{T}\) 紧且 \(\mathbf{f}\) 对拓扑 \(\sigma(\mathrm{F}, \mathrm{F}')\) 连续的情形；然后像 \(b\) 那样结束证明。）
\end{enumerate}

¶ 26) 设 S、T 是两个紧空间，f 是定义在 S × T 上的有限数值函数。

\begin{enumerate}
\def\labelenumi{\alph{enumi})}
\item
  为使映射 \(s \mapsto f(s, \cdot)\) （S 到 \(\mathbf{R}^{\mathrm{T}}\) 中的）是 S 到空间 \(\mathcal{C}(\mathrm{T})\) （赋以拓扑 \(\sigma(\mathcal{C}(\mathrm{T}), \mathcal{M}(\mathrm{T}))\) ）中的连续映射，必须且只需 \(f\) 有界，且每个部分映射 \(f(s, \cdot), f(\cdot, t)\) ( \(s \in \mathrm{S}, t \in \mathrm{T}\) ) 都连续。（为证条件充分，先证明 S 在 \(s \mapsto f(s, \cdot)\) 下的像 M 对 \(\sigma(\mathcal{C}(\mathrm{T}), \mathcal{M}(\mathrm{T}))\) 是紧的。为此，注意 \(s \mapsto f(s, \cdot)\) 当 \(\mathcal{C}(\mathrm{T})\) 赋以 T 上逐点收敛的拓扑时是连续的；利用 Eberlein 定理（TVS，IV，§5，第 3 号，定理 1）化为证明每个序列 \((f(s_n, \cdot))\) 在 \(\mathcal{C}(\mathrm{T})\) 中对 \(\sigma(\mathcal{C}(\mathrm{T}), \mathcal{M}(\mathrm{T}))\) 有一个聚点。通过考虑 T 的一个适当商空间（参见习题 10 a)）化为 T 可度量化的情形，并注意，在 \(\mathcal{C}(\mathrm{T})\) 中对 T 上逐点收敛拓扑相对紧的子集上，该拓扑与在 T 的稠密子集上逐点收敛的拓扑相同。由此得到 \((f(s_n, \cdot))\) 的一个在 T 上逐点收敛的子序列，并应用 Lebesgue 定理。最后注意，在 M 上拓扑 \(\sigma(\mathcal{C}(\mathrm{T}), \mathcal{M}(\mathrm{T}))\) 与逐点收敛的拓扑相同。）
\item
  假设每个部分映射 \(f(s, \cdot), f(\cdot, t)\) 都连续 ( \(s \in S\) , \(t \in T\) )。证明：对每个正测度 \(\mu\) ( \(S\) 上的)及每个 \(\varepsilon > 0\) ，存在紧子集 \(K \subset S\) ，使得 \(\mu(S - K) \leqslant \varepsilon\) 且 \(f\) 在 \(K \times T\) 上的限制连续。（化为 \(f\) 有界的情形；应用 a) 与习题 25 c)。）
\item
  在与 b) 相同的假设下，证明 f 对 S × T 上每个测度 \(\nu\) 都可测。（对 \(\nu\) 在 S × T 到 S 上的射影下的像应用 b)。）
\end{enumerate}

\begin{enumerate}
\def\labelenumi{\arabic{enumi})}
\setcounter{enumi}{26}
\tightlist
\item
  设 \(\mathbf{m}\) 是 \(\mathbf{T}\) 上的向量测度，取值于一个 Hausdorff 局部凸空间 \(\mathbf{F}\) 。
\end{enumerate}

\begin{enumerate}
\def\labelenumi{\alph{enumi})}
\item
  证明：若 \(\mathbf{m}\) 的支集有限，则 \(\mathbf{m} = \sum_{i=1}^{n} \mathbf{c}_i \varepsilon_{a_i}\) ，其中 \(\mathbf{c}_i \in \mathbf{F}\) 。
\item
  若 \(\mathbf{F}\) 是 Banach 空间，且 \(\mathbf{m}\) 对紧收敛拓扑连续，证明 \(\mathbf{m}\) 具有紧支集。
\item
  给出一个测度的例子：取值于 \(\mathbf{R}^{\mathbf{N}}\) ，支集非紧，且对紧收敛拓扑连续。
\end{enumerate}

\setcounter{exercisegroup}{3}
\edef\CurrentExerciseGroup{\arabic{exercisegroup}}
